\chapter[Considerações Finais]{CONSIDERAÇÕES FINAIS}

Ao longo das quatro AGES, passei de tarefas de implementação, como o
componente de entrada de dados no Decola, para decisões de arquitetura no
VanGO e o acompanhamento da equipe no Hangy. A orientação que recebi dos
colegas me ajudou, depois, a revisar soluções e explicar decisões a outros
integrantes.

Usei os conhecimentos de requisitos para entender o que deveria ser
construído e organizar as histórias de usuário. O trabalho com banco de dados,
testes e arquitetura
apareceu nas implementações e nas revisões. Na AGES II, por exemplo,
configurei os testes com ajuda de um colega de AGES IV. Na AGES III, passei
a orientar o trabalho no \textit{backend} e a cuidar de integrações e
infraestrutura. Na AGES IV, usei os conteúdos de gestão para acompanhar as
dependências, os prazos e a participação dos colegas.

Ao orientar os colegas, precisei rever minhas explicações quando eram
difíceis para quem estava começando.
No VanGO, demorei a me aproximar dos integrantes menos experientes. No Hangy,
procurei acompanhá-los mais cedo e tive de aceitar a orientação de assumir
menos implementação. Ainda considero a delegação um ponto em que preciso
melhorar, principalmente quando o prazo está próximo e vejo uma tarefa
bloqueada.

No Indústria Digital, fiz amizades que mantive ao longo do curso, mas a
entrega ficou abaixo das demais AGES. Fiquei travado nas primeiras
\textit{sprints} e precisei rever minha participação. No VanGO, a equipe
entregou um escopo grande e reservou tempo para melhorar o produto. No
Hangy, acompanhei o crescimento da autonomia dos colegas.

No início do curso, imaginava que trabalhar com Engenharia de Software
exigiria decorar a sintaxe das linguagens, pesquisar dúvidas no Stack
Overflow e passar incontáveis horas escrevendo código. Minha expectativa
se concentrava na programação.

Ao longo das quatro AGES e das experiências profissionais, passei a dar
mais atenção à compreensão das regras de negócio, à elaboração de diagramas
para organizar soluções e à comparação entre diferentes formas de escrever
o mesmo código. Hoje, essas decisões ocupam uma parte maior da minha visão
do trabalho. Também passei a delegar tarefas repetitivas à IA, reservando
mais atenção para entender o problema e avaliar a solução produzida.

Essa mudança reforçou a importância que atribuo ao domínio dos fundamentos
e ao aprofundamento técnico. Para revisar código gerado por IA, preciso
compreender seu funcionamento, reconhecer suas limitações e verificar se
atende às regras do projeto. Quero continuar aprofundando meus conhecimentos
para fazer essa avaliação com mais segurança e prevenir problemas
decorrentes do uso inadequado dessas ferramentas.

Considero que o formato da AGES foi pensado para um contexto anterior ao
uso amplo da IA na programação. Nesse formato, desenvolver uma tarefa
exigia do aluno um contato mais direto com os fundamentos, e vejo nisso
uma das grandes contribuições da disciplina à minha formação. Hoje, me
preocupa que um iniciante consiga acelerar a entrega com IA e deixe de
passar por dificuldades das quais poderia aprender. Ao mesmo tempo, o uso
dessas ferramentas faz parte do trabalho para o qual estamos nos preparando.
Gostaria que a avaliação da disciplina acompanhasse essa mudança.

Acredito que vale a pena implementar por conta própria, pelo menos na
primeira experiência com cada tipo de tarefa. Vejo uma semelhança com os
exercícios de matemática feitos sem calculadora para desenvolver o
raciocínio. Ao escrever uma solução, o estudante precisa lidar com dúvidas
que podem passar despercebidas quando recebe o código pronto. Depois dessa
primeira experiência, terá uma base melhor para avaliar as sugestões da IA
e reconhecer quando precisa investigá-las.

Fiscalizar o uso dessas ferramentas seria difícil. Minha sugestão é dar
mais peso à demonstração do aprendizado na avaliação da AGES. Proponho que cada
\ac{pr} integrado seja acompanhado de um vídeo em que o autor explique
a solução, incluindo as partes geradas com IA. A explicação
deveria mostrar o funcionamento do código e as decisões tomadas. Ao preparar o
vídeo, o aluno precisaria estudar o que está entregando e
poderia identificar dúvidas que ainda não resolveu.

Os vídeos também serviriam para consultas posteriores. Por mais de um ano,
voltei ao Decola quando tinha dúvidas de \textit{frontend}. Explicações em
vídeo teriam me ajudado a entender aquelas implementações. Gostaria de
preservar na AGES essa possibilidade de aprender com um projeto real e
voltar a ele depois.
